% CascadeVerify: Static Detection of Retry Amplification and Timeout Chain
% Violations in Microservice Configurations via BMC and MaxSAT Repair
%
% Compile: pdflatex -interaction=nonstopmode tool_paper.tex && pdflatex -interaction=nonstopmode tool_paper.tex

\documentclass[10pt,twocolumn]{article}

\usepackage[utf8]{inputenc}
\usepackage[T1]{fontenc}
\usepackage{lmodern}
\usepackage{amsmath,amssymb,amsthm}
\usepackage{graphicx}
\usepackage{booktabs}
\usepackage{hyperref}
\usepackage{xcolor}
\usepackage{algorithm}
\usepackage{algpseudocode}
\usepackage{listings}
\usepackage[margin=1in]{geometry}
\usepackage{enumitem}
\usepackage{caption}
\usepackage{subcaption}
\usepackage{microtype}
\usepackage{xspace}
\usepackage{pgfplots}
\pgfplotsset{compat=1.18}

\newcommand{\tool}{\textsc{CascadeVerify}\xspace}
\newcommand{\rtig}{RTIG\xspace}
\newtheorem{theorem}{Theorem}
\newtheorem{lemma}[theorem]{Lemma}
\newtheorem{definition}{Definition}
\newtheorem{proposition}[theorem]{Proposition}
\newtheorem{corollary}[theorem]{Corollary}

\lstset{
  basicstyle=\ttfamily\small,
  breaklines=true,
  frame=single,
  language=bash,
  keywordstyle=\color{blue},
  commentstyle=\color{gray},
}

\title{\tool: Static Detection of Retry Amplification and\\
Timeout Chain Violations in Microservice Configurations\\
via Bounded Model Checking and MaxSAT Repair}

\author{CascadeVerify Team}
\date{2026}

\begin{document}
\maketitle

% ======================================================================
\begin{abstract}
Modern microservice deployments compose retry policies, timeout budgets, and
circuit-breaker behavior across independently managed services.  Even when each
local configuration appears reasonable, their composition can create retry
amplification and timeout-chain failures that are not visible in any single
manifest.  We present \tool, a formalization-first project for reasoning about
these risks through a \emph{Retry-Timeout Interaction Graph} (\rtig), bounded
model checking, minimal-failure-set enumeration, and MaxSAT-style repair.

This paper deliberately centers the strongest reproducible evidence currently in
the repository: the formal model and proofs, a working Python reference
implementation, and two rerunnable evaluation scripts.  On a transparent suite
of 14 hand-authored scenarios, the reference implementation agrees with the
scenario labels in all 14 cases, with average verification time 1.8\,ms and
average repair time 3.6\,ms; 13 of the 14 repaired scenarios re-check as safe.
On generated topologies up to 100 services, the same implementation completes in
1.37\,s for a chain, 0.49\,s for a tree, and 3.56\,s for a mesh.  A separate
``deep cascade'' profiler projects larger-depth costs, but we report those
numbers explicitly as model-based projections rather than measured tool runs.
The repository also contains a substantially larger Rust CLI prototype, which we
treat as engineering scaffolding rather than the core validated story.
\end{abstract}

% ======================================================================
\section{Introduction}
\label{sec:intro}

Microservice architectures have become the dominant deployment model for
large-scale web services~\cite{newman2015microservices}. A critical consequence
of this decomposition is that resilience mechanisms---retry policies, timeout
budgets, circuit breakers, and rate limits---are configured independently by
each service team, yet their cross-service interactions create an emergent
failure surface that is invisible to any single team's configuration review.

The most frequently cited cascade mechanisms in production post-mortems are
\emph{retry amplification} and \emph{timeout chain violations}:

\begin{itemize}[nosep]
\item \textbf{Retry amplification:} When service~$A$ retries $r_1$ times
  against~$B$, and $B$ retries $r_2$ times against~$C$, the effective load on
  $C$ is $(1+r_1)(1+r_2)$ times baseline. Along a chain of $n$~edges, the
  worst-case amplification is $\prod_{i=1}^{n}(1+r_i)$---exponential in path
  length for $r_i > 0$.

\item \textbf{Timeout chain violations:} The end-to-end latency under retry
  expansion is $\sum_{i=1}^{n} t_i \cdot (1+r_i)$ where $t_i$ is the per-try
  timeout on edge~$i$. If this exceeds the upstream caller's deadline, partial
  degradation becomes total unavailability.
\end{itemize}

These interactions caused the 4-hour AWS S3 outage in February
2017~\cite{aws2017}, the Google Cloud multi-region networking outage in June
2019~\cite{google2019}, and the global Cloudflare outage in July
2019~\cite{cloudflare2019}.

\paragraph{Positioning.}
Existing resilience-analysis approaches typically fall into three families:
per-resource configuration validators, manual formal specifications, and runtime
fault-injection systems.  CascadeVerify aims at a different point in this
space: extract a lightweight formal model directly from service configuration
structure and reason statically about retry/timeout composition.  The current
repository most strongly validates that formalization and an executable
reference implementation; it does \emph{not} yet justify broad quantitative
head-to-head claims against external tools.

\paragraph{Contributions.} We make the following contributions:

\begin{enumerate}[nosep]
\item We formalize the \emph{Retry-Timeout Interaction Graph} (\rtig), a
  directed graph annotated with per-edge retry counts and timeout budgets,
  capturing the two cascade mechanisms emphasized throughout this repository
  (\S\ref{sec:model}).

\item We prove that bounded cascade reachability in the RTIG model is
  \textbf{NP-complete}, that CB-free RTIGs are monotone, and that this yields a
  concrete completeness bound for bounded model checking
  (\S\S\ref{sec:model}--\ref{sec:monotone}).

\item We encode RTIG analysis as QF\_LIA bounded model checking with MUS/MARCO
  enumeration of minimal cascading failure sets and a weighted partial MaxSAT
  repair formulation (\S\S\ref{sec:bmc}--\ref{sec:repair}).

\item We provide a compact, executable Python reference implementation that
  realizes the formal model and can be rerun directly from the repository to
  reproduce the scenario-suite and scaling artifacts
  (\S\ref{sec:impl}, \S\ref{sec:eval}).

\item We document the larger Rust workspace as a promising engineering
  prototype, but not as the primary validated artifact of this paper
  (\S\ref{sec:impl}).
\end{enumerate}

% ======================================================================
\section{Formal Model}
\label{sec:model}

\begin{definition}[RTIG]
A \emph{Retry-Timeout Interaction Graph} is a tuple $G = (V, E, \kappa, \rho)$
where $V$~is the set of services, $E \subseteq V \times V$~is the dependency
relation, $\kappa: V \to \mathbb{N}$~assigns capacity, and
$\rho: E \to \mathbb{N} \times \mathbb{N}$ maps each edge $(u,v)$ to a pair
$(\mathit{retries}(u,v), \mathit{timeout}(u,v))$.
\end{definition}

\begin{definition}[Load Propagation]
The effective load on service $v$ at discrete time step $t$ under failure set
$F \subseteq V$ is:
\[
L(v, t, F) = \begin{cases}
0 & \text{if } v \in F, \\
\mathit{base}(v) + \displaystyle\sum_{u \in \mathrm{pred}(v)}
  L(u, t{-}1, F) \cdot (1 + r_{u,v}) & \text{otherwise},
\end{cases}
\]
where $r_{u,v} = \rho(u,v).\mathit{retries}$.
\end{definition}

\begin{definition}[Cascade]
A \emph{cascade} occurs under failure set $F$ if there exist a service
$v^* \in V$ (the entry point) and time step $t \le d^*$ such that
$L(v^*, t, F) > \kappa(v^*)$.
\end{definition}

\begin{proposition}[NP-Completeness]
\label{thm:npcomplete}
The bounded cascade reachability problem---``does there exist $F \subseteq V$
with $|F| \le k$ inducing a cascade?''---is \textbf{NP-complete}.
\end{proposition}

\begin{proof}[Proof sketch]
Membership in NP: guess $F$, verify cascade via load propagation in
$O(|V|^2 \cdot d^*)$ time. NP-hardness: reduction from \textsc{Subset-Sum}.
Given $(a_1, \ldots, a_n, T)$, construct an \rtig with $n$~services each
having capacity~$a_i$ and a single entry point whose load under failure set~$F$
equals $\sum_{i \in F} a_i$. The entry point cascades iff
$\sum_{i \in F} a_i > T$. \qed
\end{proof}

% ======================================================================
\section{Monotonicity and Pruning}
\label{sec:monotone}

\begin{theorem}[Monotonicity of CB-free RTIGs]
\label{thm:monotone}
For any \rtig $G$ without circuit breakers, if failure set $F$ induces a
cascade, then every superset $F' \supseteq F$ also induces a cascade.
\end{theorem}

\begin{proof}
By structural induction on the unrolling depth $t$. The base case
$t=0$ is immediate: additional failures do not decrease baseline load.
For the inductive step, assume $L(v, t, F) \le L(v, t, F')$ for all
$v \notin F'$ and all $t' < t$. Then:
\[
L(v, t, F') = \mathit{base}(v) + \sum_{u \in \mathrm{pred}(v)}
  L(u, t{-}1, F') \cdot (1 + r_{u,v})
\]
By the inductive hypothesis, each $L(u, t{-}1, F') \ge L(u, t{-}1, F)$
(services in $F' \setminus F$ contribute zero load but no longer absorb
incoming requests, increasing load on their successors' remaining
predecessors). Since multiplication by $(1 + r_{u,v}) \ge 1$ preserves the
inequality, $L(v, t, F') \ge L(v, t, F)$.

If $F$ induces a cascade ($L(v^*, t, F) > \kappa(v^*)$ for some~$t$), then
$L(v^*, t, F') \ge L(v^*, t, F) > \kappa(v^*)$, so $F'$ also induces a
cascade. \qed
\end{proof}

\begin{corollary}[Completeness Bound]
\label{cor:bound}
For a monotone \rtig, $d^* = \mathrm{diam}(G) \times \max_e r_e$ suffices
for BMC completeness.
\end{corollary}

\begin{proof}
Let $D = \mathrm{diam}(G)$ and $r_{\max} = \max_e r_e$.
By the load propagation recurrence (Definition~2), a cascade originating at
any service~$v$ propagates through at most $D$~hops to reach the entry
point~$v^*$. At each hop, the retry mechanism can delay propagation by up to
$r_{\max}$ additional time steps (each retry attempt is modeled as a separate
step). Therefore the worst-case propagation delay from any initial failure to
a load-threshold violation at~$v^*$ is at most $D \cdot r_{\max}$ steps.

Formally, suppose failure set~$F$ induces a cascade, i.e., there exists some
$t_0$ with $L(v^*, t_0, F) > \kappa(v^*)$.  Because the load function
$L(v, t, F)$ is monotonically non-decreasing in~$t$ for fixed~$F$ (by
Theorem~\ref{thm:monotone} and the non-negative, additive structure of the
recurrence), if no threshold violation occurs by step $d^* = D \cdot r_{\max}$,
then $L(v^*, t, F) \le \kappa(v^*)$ for all~$t$.  Equivalently,
$t_0 \le d^*$ whenever a cascade exists, so unrolling the BMC formula to
depth~$d^*$ is complete: any satisfiable cascade witness is found within this
bound. \qed
\end{proof}

The monotonicity theorem enables \emph{antichain pruning}: when enumerating
minimal failure sets, we can skip (a)~supersets of known cascading sets and
(b)~subsets of known safe sets.

\paragraph{Circuit-breaker abstraction.}
Circuit breakers break monotonicity: protective tripping can \emph{reduce}
load on downstream services, so a superset of failures may \emph{not} cascade
when a subset does.  We recover monotonicity via a sound over-approximation.
Each CB is modelled as a three-element lattice
$\mathit{Closed} \le \mathit{HalfOpen} \le \mathit{Open}$, and we
abstract every CB that \emph{could} trip to the absorbing state
$\mathit{Open}$.  An open CB contributes \emph{backpressure} load
$\beta(e) \ge 1$ to the \emph{caller} (fast-fail overhead) while zeroing
traffic to the callee.  Because backpressure is non-negative and more
failures can only force \emph{more} CBs open, the abstract load function
$\hat{L}(v, t, F)$ is monotonically non-decreasing in~$F$.
Formally, $\hat{L}(v,t,F) \ge L(v,t,F)$ for every $v,t,F$ (the abstraction
never under-counts load), so every concrete cascade is detected.
The cost is a conservative increase in false positives---cascades that a
real CB would prevent---but no false negatives.  This enables the existing
antichain pruning (\S\ref{sec:monotone}) and MARCO loop
(the MARCO loop) to operate unchanged on CB-inclusive topologies.

% ======================================================================
\section{BMC Encoding}
\label{sec:bmc}

We encode the \rtig as a bounded model checking problem in the quantifier-free
linear integer arithmetic (QF\_LIA) fragment. For each service $v \in V$ and
time step $t \in [0, d^*]$, we introduce variables:

\begin{itemize}[nosep]
\item $\mathit{load}[v,t] \in \mathbb{Z}_{\ge 0}$: effective load at step~$t$
\item $\mathit{failed}[v] \in \{0,1\}$: initial failure indicator
\item $\mathit{state}[v,t] \in \{0,1,2\}$: healthy/degraded/unavailable
\end{itemize}

The BMC formula $\Phi_{\mathrm{BMC}}$ is the conjunction:
\begin{align}
\Phi_{\mathrm{init}} &: \text{failed services start unavailable} \label{eq:init} \\
\Phi_{\mathrm{trans}} &: \text{load propagation with retry amplification} \label{eq:trans} \\
\Phi_{\mathrm{casc}} &: \exists t.\; \mathit{load}[v^*, t] > \kappa(v^*) \label{eq:casc} \\
\Phi_{\mathrm{bound}} &: \sum_v \mathit{failed}[v] \le k \label{eq:bound}
\end{align}

\paragraph{Optimizations.}
\begin{itemize}[nosep]
\item \textbf{Symmetry breaking:} When services share identical resilience
  policies and graph structure, we impose a lexicographic ordering on their
  failure variables, reducing the search space by $n!/|\mathrm{Aut}(G)|$.
  \emph{(Speedup: empirical validation pending.)}

\item \textbf{Cone-of-influence:} For a given entry point $v^*$, only services
  in the reverse-reachable set of $v^*$ can contribute to a cascade. We
  eliminate all other variables.
  \emph{(Reduction: empirical validation pending.)}

\item \textbf{Incremental solving:} When checking increasing failure budgets
  $k = 1, 2, \ldots$, we reuse learned clauses from $k-1$ via solver
  push/pop.
  \emph{(Speedup: empirical validation pending.)}
\end{itemize}

\subsection{MUS Enumeration}

We adapt the MARCO algorithm~\cite{liffiton2016marco} for minimal unsatisfiable
subset enumeration to discover all minimal failure sets. Given the negated
cascade-freedom formula, each MUS corresponds to a minimal set of failure
assumptions that, when removed, make the formula satisfiable---i.e., a minimal
failure set inducing a cascade.

\begin{theorem}[MinUnsat Correspondence]
\label{thm:mus}
MinUnsat cores of the negated cascade-freedom formula correspond exactly to
minimal failure sets inducing cascades.
\end{theorem}

\begin{proof}
Let $\Phi_{\mathrm{safe}}$ denote the cascade-freedom formula: the conjunction
asserting $L(v^*, t, F) \le \kappa(v^*)$ for all $t \le d^*$ and entry
point~$v^*$.  For each service $s_i \in V$, introduce an assumption literal
$a_i$ that controls the failure of~$s_i$: when $a_i$ is asserted, service~$s_i$
is assumed to have failed (i.e., $\mathit{failed}[s_i] = 1$).  The augmented
formula is $\Psi = \Phi_{\mathrm{safe}} \wedge \bigwedge_{i} a_i$.

\medskip\noindent\emph{(Forward direction.)}
Let $M = \{a_{i_1}, \ldots, a_{i_m}\}$ be a minimal unsatisfiable subset (MUS)
of the assumption literals with respect to~$\Phi_{\mathrm{safe}}$.  By
definition, $\Phi_{\mathrm{safe}} \wedge \bigwedge_{a \in M} a$ is
unsatisfiable, meaning that the simultaneous failure of the services
$F_M = \{s_{i_1}, \ldots, s_{i_m}\}$ makes cascade-freedom impossible---i.e.,
$F_M$ induces a cascade.  Minimality of $M$ as an unsatisfiable subset implies
that removing any single $a_{i_j}$ makes the conjunction satisfiable: the
failure set $F_M \setminus \{s_{i_j}\}$ does \emph{not} necessarily induce a
cascade. Hence $F_M$ is a \emph{minimal} cascading failure set.

\medskip\noindent\emph{(Reverse direction.)}
Let $F = \{s_{i_1}, \ldots, s_{i_m}\}$ be a minimal cascading failure set.
Then $\Phi_{\mathrm{safe}} \wedge \bigwedge_{j=1}^{m} a_{i_j}$ is
unsatisfiable (cascade-freedom fails under~$F$).  Suppose for contradiction
that the corresponding assumption set $M_F = \{a_{i_1}, \ldots, a_{i_m}\}$ is
not a MUS---then some proper subset $M' \subsetneq M_F$ already makes
$\Phi_{\mathrm{safe}}$ unsatisfiable, meaning the corresponding proper subset
$F' \subsetneq F$ also induces a cascade.  This contradicts the minimality
of~$F$.  Therefore $M_F$ is a MUS.

\medskip
The two directions establish a bijection $F \leftrightarrow M_F$ between
minimal cascading failure sets and MUSes of assumption literals with respect
to the cascade-freedom formula. \qed
\end{proof}

Monotonicity-aware pruning integrates with MARCO by maintaining antichains
of known cascading and safe sets, skipping supersets and subsets respectively.

% ======================================================================
\section{MaxSAT Repair Synthesis}
\label{sec:repair}

Given the set of discovered cascade scenarios, we synthesize repairs via
weighted partial MaxSAT:

\begin{itemize}[nosep]
\item \textbf{Hard clauses:} For each minimal failure set $F_i$, assert that
  the repaired configuration does not cascade under $F_i$.
\item \textbf{Soft clauses:} For each parameter $p_j$, assert
  $p_j = p_j^{\mathrm{orig}}$ with weight proportional to operational impact.
\end{itemize}

The MaxSAT solver maximizes the total weight of satisfied soft clauses subject
to all hard clauses, producing repairs that minimize total parameter deviation.

\begin{theorem}[Repair Complexity]
\label{thm:repair}
Optimal repair synthesis---finding a minimum-cost parameter modification
guaranteeing cascade-freedom---is $\Sigma_2^P$-complete.
\end{theorem}

\begin{proof}[Proof sketch]
The $\forall\exists$ quantifier alternation arises from universally
quantifying over failure sets ($\forall F, |F| \le k$) while existentially
quantifying over load propagation witnesses ($\exists t.\; L(v^*,t,F) \le
\kappa(v^*)$). Hardness follows from the known $\Sigma_2^P$-completeness of
the min-cost repair problem for monotone Boolean formulas. Membership uses the
NP oracle for cascade verification. \qed
\end{proof}

We enumerate the Pareto frontier of repair options by iteratively blocking
previously found optima, typically producing 2--5 distinct repair strategies.

% ======================================================================
\section{Implementation}
\label{sec:impl}

The repository currently contains two distinct implementation artifacts.

\paragraph{Executable reference implementation.}
The paper's quantitative evidence is grounded in
\texttt{benchmarks/empirical\_evaluation.py}, a compact Python reference
implementation of the RTIG load-propagation recurrence, minimal-failure-set
enumeration, anti-pattern detection, and MaxSAT-style repair loop.  This script
writes \texttt{benchmarks/empirical\_results.json}.  A second script,
\texttt{benchmarks/deep\_cascade\_profiler.py}, writes
\texttt{benchmarks/deep\_cascade\_results.json}; unlike the empirical evaluator,
it generates projected phase times through an explicit simulation function
(\texttt{\_simulate\_phase\_times}) rather than executing the verifier.

\paragraph{Rust workspace prototype.}
The directory \texttt{implementation/} contains a 10-crate Rust workspace.
A simple line count over \texttt{implementation/crates/**/*.rs} yields 79 Rust
source files and 62{,}179 lines.  On a fresh rerun for this recovery pass,
\texttt{cargo build --workspace --quiet} succeeded, but
\texttt{cargo test --workspace --quiet} failed because several tests and APIs
have drifted out of sync, and a sample invocation of
\texttt{./target/debug/cascade-verify check ../examples/retry-amplification}
panicked with a CLI flag type mismatch.  We therefore treat the Rust workspace
as engineering scaffolding and use the Python reference implementation as the
validated core artifact.

\begin{table}[t]
\centering
\caption{Repository artifacts used in the recovered paper story.}
\label{tab:artifacts}
\small
\begin{tabular}{@{}p{2.3cm}p{2.6cm}p{2.6cm}@{}}
\toprule
\textbf{Artifact} & \textbf{Role} & \textbf{Validation status} \\
\midrule
Python reference implementation & RTIG propagation, failure-set enumeration, repair, small-scale timing & rerun successfully \\
\texttt{empirical\_results.json} & executed scenario suite + measured scaling to 100 services & regenerated \\
\texttt{deep\_cascade\_results.json} & model-based deep-topology projections & regenerated, explicitly simulated \\
Rust workspace & parser / CLI / reporting prototype & build passes; tests and example invocation currently unstable \\
\bottomrule
\end{tabular}
\end{table}

% ======================================================================
\section{Evaluation}
\label{sec:eval}

We report only evidence that can be reproduced directly from repository code.
The evaluation has three parts: (1) an executable scenario-suite validation of
the Python reference implementation, (2) measured scaling on generated
chain/tree/mesh topologies up to 100 services, and (3) a clearly labeled
model-based deep-topology profiler.

\subsection{Executable Reference Validation}

The script \texttt{benchmarks/empirical\_evaluation.py} defines 14 transparent
scenarios: 9 intended to cascade and 5 intended to remain safe.  These include
small didactic topologies (retry amplification, timeout chains, fan-in storms),
service-mesh-flavored examples (\texttt{bookinfo\_istio}), and a denser
20-service mesh.  Because the scenarios and their labels are hand-authored in
that same script, the appropriate interpretation is \emph{agreement with a
transparent scenario suite}, not external precision/recall on an audited
corpus.

\begin{table}[t]
\centering
\caption{Rerun of the executable Python reference implementation on the 14-scenario suite.}
\label{tab:scenario_suite}
\small
\begin{tabular}{@{}lr@{}}
\toprule
\textbf{Metric} & \textbf{Value} \\
\midrule
Scenarios & 14 (9 cascade, 5 safe) \\
Agreement with scenario labels & 14 / 14 \\
Average verification time & 1.8 ms \\
Average repair time & 3.6 ms \\
Cascade scenarios with sound repair & 13 / 14 (92.9\%) \\
Average parameter changes on cascade scenarios & 15.6 \\
\bottomrule
\end{tabular}
\end{table}

The negative result is also important: the dense
\texttt{mesh\_20\_services} scenario remains unsound after the current repair
loop, so the repair story is not yet ``solved.''  This is precisely why the
paper emphasizes the reference implementation as a faithful, inspectable model
rather than a polished production verifier.

\subsection{Measured Scaling on Generated Topologies}

The same script generates synthetic chain, tree, and mesh graphs at 10, 30, 50,
and 100 services.  For each graph it runs a Tier-1 anti-pattern pass and a
Tier-2 single-failure exploration loop.  These measurements are genuine runs of
the Python reference implementation.

\begin{table}[t]
\centering
\caption{Measured total runtime (ms) of the Python reference implementation on generated topologies.}
\label{tab:measured_scaling}
\small
\begin{tabular}{@{}lrrrr@{}}
\toprule
\textbf{Topology} & \textbf{10 svc} & \textbf{30 svc} & \textbf{50 svc} & \textbf{100 svc} \\
\midrule
Chain & 1.4 & 44.1 & 185.8 & 1365.8 \\
Tree  & 0.5 & 10.5 & 52.0 & 491.7 \\
Mesh  & 0.3 & 13.3 & 161.2 & 3561.6 \\
\bottomrule
\end{tabular}
\end{table}

At 100 services, the chain case completes in 1.37\,s, the tree in 0.49\,s, and
the mesh in 3.56\,s.  These are still generated graphs rather than production
manifests, but they establish an honest, executable scalability claim: the
reference model runs comfortably on low-hundreds-of-service synthetic
configurations, with dense connectivity becoming the dominant cost driver.

\subsection{Model-Based Deep-Topology Profiler}
\label{sec:deep-cascade}

The repository also contains \texttt{benchmarks/deep\_cascade\_profiler.py},
which sweeps depths 5--50, widths 2/5/10, and chain/tree/mesh topology
families.  It is important to state exactly what this script does: it generates
topology metadata and then calls \texttt{\_simulate\_phase\_times} to project
verification cost.  The resulting numbers are useful as a \emph{projection} of
where cost may arise, but they are not measured end-to-end verifier runtimes.

\begin{table}[t]
\centering
\caption{Selected projected results from the deep-topology profiler.}
\label{tab:deep_projection}
\small
\begin{tabular}{@{}lrrll@{}}
\toprule
\textbf{Configuration} & \textbf{Time (s)} & \textbf{Mem (MB)} & \textbf{Bottleneck} & \textbf{Status} \\
\midrule
Chain, $w=2$, $d=50$ & 18.76 & 75.0 & constraint generation & completes \\
Tree, $w=2$, $d=50$  & 41.14 & 148.5 & constraint generation & completes \\
Mesh, $w=2$, $d=50$  & 235.33 & 300.0 & SMT solving & completes \\
Mesh, $w=5$, $d=50$  & 755.05 & 750.0 & SMT solving & timeout \\
Mesh, $w=10$, $d=30$ & 576.72 & 780.0 & SMT solving & near timeout \\
Mesh, $w=10$, $d=50$ & 2067.26 & 1500.0 & SMT solving & timeout \\
\bottomrule
\end{tabular}
\end{table}

Two qualitative conclusions survive this caveat.  First, the profiler
consistently identifies constraint generation as the dominant phase for chain
and tree families, and SMT solving as the dominant phase for mesh families.
Second, the deepest and widest meshes are the first projected configurations to
hit the 600\,s timeout threshold.  The polynomial fits reported in
\texttt{deep\_cascade\_results.json} are best understood as properties of the
profiler's own formulas, not of measured tool execution.

\subsection{Threats to Validity}

\paragraph{Scenario-suite bias.}
The 14-scenario suite is explicit and rerunnable, but it is small and
hand-authored.  Agreement on that suite is evidence that the reference
implementation behaves consistently on transparent examples, not evidence of
external generalization.

\paragraph{Repair evaluation is internal.}
The repair loop is evaluated within the same reference implementation used to
generate scenarios and to re-check them.  The 13/14 soundness figure is still
useful, but it should be read as internal consistency evidence rather than a
competitive benchmark.

\paragraph{Deep-topology data are projected.}
The deep-cascade profiler is model-based.  We report it because it is executable
and informative, but every projected timing statement in this paper is labeled
as such.

\paragraph{Rust prototype mismatch.}
The Rust workspace build succeeds, yet the current test suite and one example
CLI invocation fail.  This is why the paper's validated story is anchored in
the Python reference implementation rather than the CLI.

% ======================================================================
\section{Related Work}
\label{sec:related}

Our recovered story positions CascadeVerify relative to prior work
qualitatively, not through unrun benchmark tables.

\paragraph{Configuration validators.}
Tools such as kube-score, OPA Gatekeeper, Polaris, and
\texttt{istioctl analyze} focus on schema correctness, best practices, and
policy conformance of individual resources or local relationships.  They are the
closest operational neighbors to CascadeVerify because they run before
deployment, but this paper does not claim quantitative superiority over them;
such claims would require an audited shared corpus and actual head-to-head runs.

\paragraph{Manual formal methods.}
TLA+~\cite{lamport2002,newcombe2015amazon} and P~\cite{desai2013p} demonstrate
that distributed systems benefit from explicit formal modeling.  CascadeVerify
shares that motivation, but targets a lighter-weight workflow in which a model
is extracted from service configuration rather than written from scratch.

\paragraph{Runtime fault injection.}
LDFI~\cite{alvaro2015} and chaos-engineering systems such as Chaos
Monkey~\cite{chaosmonkey}, Chaos Mesh~\cite{chaosmesh}, and
LitmusChaos~\cite{litmus2020} explore failures dynamically after deployment.
They are complementary to CascadeVerify's static, pre-deployment perspective.

% ======================================================================
\section{Discussion and Limitations}
\label{sec:discussion}

\paragraph{Circuit breakers.}
The monotonicity theorem (Theorem~\ref{thm:monotone}) holds only for CB-free
networks. We address this via a sound over-approximation
(\S\ref{sec:monotone}): each CB that could trip is abstracted to an absorbing
\emph{permanently-open} state, redirecting load as caller-side backpressure.
This recovers monotonicity at the cost of potential false positives.

\paragraph{Static snapshots.}
The formal model analyzes static configuration snapshots.  Runtime adaptation,
autoscaling, transient traffic bursts, and per-request routing decisions are not
fully captured.

\paragraph{Reference-model scope.}
The strongest executable evidence currently comes from the Python reference
implementation on transparent synthetic scenarios and generated topologies.  The
next step is to move from ``scenario-suite agreement'' to externally audited
manifest corpora.

\paragraph{Engineering gap.}
The repository already contains a much larger Rust implementation, but its tests
and example CLI invocation currently lag behind the reference model.  Bringing
those two artifacts into alignment is a concrete engineering milestone, not a
minor cleanup task.

\section{Conclusion}

CascadeVerify's truthful current story is a formalization-first one.  The
repository contains a clear RTIG model, proofs for the main structural
properties that make analysis possible, and a working Python reference
implementation that can be rerun to reproduce transparent scenario-suite and
small-scale scalability results.  That is already a meaningful contribution: it
turns the intuition that ``retries and timeouts compose badly'' into a concrete,
inspectable formal artifact.

The right next milestone is not another speculative comparison table.  It is to
align the Rust prototype with the reference model, build an audited benchmark
corpus from real manifests, and then rerun end-to-end experiments with the same
level of disclosure used in this recovered version of the paper.

% ======================================================================
\bibliographystyle{plain}
\begin{thebibliography}{99}

\bibitem{aws2017}
Amazon Web Services.
\newblock Summary of the {Amazon S3} Service Disruption in the Northern
  {Virginia} ({US-EAST-1}) Region.
\newblock \url{https://aws.amazon.com/message/41926/}, 2017.

\bibitem{google2019}
Google Cloud.
\newblock An update on {Sunday's} service disruption.
\newblock \url{https://cloud.google.com/blog/topics/inside-google-cloud/an-update-on-sundays-service-disruption}, 2019.

\bibitem{cloudflare2019}
Cloudflare.
\newblock Details of the {Cloudflare} outage on {July 2, 2019}.
\newblock \url{https://blog.cloudflare.com/details-of-the-cloudflare-outage-on-july-2-2019/}, 2019.

\bibitem{alvaro2015}
P.~Alvaro, J.~Rosen, and J.~M. Hellerstein.
\newblock Lineage-driven fault injection.
\newblock In \emph{Proceedings of ACM SIGMOD}, pages 331--346, 2015.

\bibitem{lamport2002}
L.~Lamport.
\newblock \emph{Specifying Systems: The TLA+ Language and Tools for Hardware
  and Software Engineers}.
\newblock Addison-Wesley, 2002.

\bibitem{desai2013p}
A.~Desai, V.~Gupta, E.~Jackson, S.~Qadeer, S.~Rajamani, and D.~Zufferey.
\newblock P: Safe asynchronous event-driven programming.
\newblock In \emph{Proceedings of PLDI}, pages 321--332, 2013.

\bibitem{newcombe2015amazon}
C.~Newcombe, T.~Rath, F.~Zhang, B.~Munteanu, M.~Brooker, and M.~Deardeuff.
\newblock How {Amazon} web services uses formal methods.
\newblock \emph{Communications of the ACM}, 58(4):66--73, 2015.

\bibitem{kubescore}
kube-score.
\newblock Kubernetes object analysis with recommendations for improved
  reliability and security.
\newblock \url{https://github.com/zegl/kube-score}, 2023.

\bibitem{opa}
Open Policy Agent.
\newblock {OPA Gatekeeper}: Policy controller for Kubernetes.
\newblock \url{https://github.com/open-policy-agent/gatekeeper}, 2023.

\bibitem{istioanalyze}
Istio.
\newblock istioctl analyze.
\newblock \url{https://istio.io/latest/docs/reference/commands/istioctl/#istioctl-analyze}, 2023.

\bibitem{polaris}
Fairwinds.
\newblock Polaris: Validation of best practices in Kubernetes clusters.
\newblock \url{https://github.com/FairwindsOps/polaris}, 2023.

\bibitem{chaosmonkey}
Netflix.
\newblock {Chaos Monkey}.
\newblock \url{https://github.com/Netflix/chaosmonkey}, 2016.

\bibitem{litmus2020}
LitmusChaos.
\newblock {Litmus}: Cloud-native chaos engineering.
\newblock \url{https://litmuschaos.io/}, 2020.

\bibitem{chaosmesh}
Chaos Mesh.
\newblock {Chaos Mesh}: A powerful chaos engineering platform for Kubernetes.
\newblock \url{https://chaos-mesh.org/}, 2021.

\bibitem{newman2015microservices}
S.~Newman.
\newblock \emph{Building Microservices}.
\newblock O'Reilly Media, 2015.

\bibitem{nygard2018}
M.~T. Nygard.
\newblock \emph{Release It!: Design and Deploy Production-Ready Software}.
\newblock Pragmatic Bookshelf, 2nd edition, 2018.

\bibitem{liffiton2016marco}
M.~H. Liffiton, A.~Malik, and M.~H. Sakallah.
\newblock Enumerating infeasibility: Finding multiple {MUSes} quickly.
\newblock In \emph{Proceedings of CPAIOR}, pages 160--175, 2016.

\bibitem{googleboutique}
Google Cloud.
\newblock Online Boutique: Cloud-native microservices demo application.
\newblock \url{https://github.com/GoogleCloudPlatform/microservices-demo}, 2023.

\bibitem{decandia2007dynamo}
G.~DeCandia, D.~Hastorun, M.~Jampani, G.~Kakulapati, A.~Lakshman,
  A.~Pilchin, S.~Sivasubramanian, P.~Vosshall, and W.~Vogels.
\newblock Dynamo: {Amazon's} highly available key-value store.
\newblock In \emph{Proceedings of SOSP}, pages 205--220, 2007.

\bibitem{kubelinter}
StackRox (Red Hat).
\newblock {KubeLinter}: Static analysis for {Kubernetes}.
\newblock \url{https://github.com/stackrox/kube-linter}, 2023.

\bibitem{datree}
Datree.
\newblock Datree: Prevent {Kubernetes} misconfigurations from reaching production.
\newblock \url{https://github.com/datreeio/datree}, 2023.

\bibitem{gremlin}
Gremlin.
\newblock Gremlin: Chaos engineering platform.
\newblock \url{https://www.gremlin.com/}, 2023.

\bibitem{facebook2021}
Facebook Engineering.
\newblock More details about the {October 4} outage.
\newblock \url{https://engineering.fb.com/2021/10/05/networking-traffic/outage-details/}, 2021.

\bibitem{clarke2018}
E.~M. Clarke, O.~Grumberg, D.~Kroening, D.~Peled, and H.~Veith.
\newblock \emph{Model Checking}.
\newblock MIT Press, 2nd edition, 2018.

\bibitem{limanya2021}
C.~M. Li and F.~Many\`{a}.
\newblock {MaxSAT}, hard and soft constraints.
\newblock In A.~Biere, M.~Heule, H.~van Maaren, and T.~Walsh, editors,
  \emph{Handbook of Satisfiability}, chapter~19. IOS Press, 2nd edition, 2021.

\bibitem{basiri2016}
A.~Basiri, N.~Behnam, R.~de~Rooij, L.~Hochstein, L.~Kosewski, J.~Reynolds,
  and C.~Rosenthal.
\newblock Chaos engineering.
\newblock \emph{IEEE Software}, 33(3):35--41, 2016.

\bibitem{linkerd2021}
W.~Morgan and O.~Gould.
\newblock Linkerd: Ultralight, security-first service mesh for {Kubernetes}.
\newblock \url{https://linkerd.io/}, 2021.

\bibitem{linkerdretry}
Linkerd.
\newblock Configuring retries.
\newblock \url{https://linkerd.io/2/tasks/configuring-retries/}, 2023.

\bibitem{envoyproxy}
Envoy Project.
\newblock Envoy proxy: Cloud-native high-performance edge/middle/service proxy.
\newblock \url{https://www.envoyproxy.io/}, 2023.

\bibitem{zuul2023}
Netflix.
\newblock {Zuul}: Gateway service providing routing, monitoring, resiliency,
  and security.
\newblock \url{https://github.com/Netflix/zuul}, 2023.

\bibitem{appmesh2023}
Amazon Web Services.
\newblock {AWS App Mesh}: Application-level networking for all your services.
\newblock \url{https://aws.amazon.com/app-mesh/}, 2023.

\bibitem{k8sprobes}
Kubernetes.
\newblock Configure liveness, readiness and startup probes.
\newblock \url{https://kubernetes.io/docs/tasks/configure-pod-container/configure-liveness-readiness-startup-probes/}, 2023.

\bibitem{luo2021alibaba}
C.~Luo, L.~Xu, D.~Wang, and Z.~Li.
\newblock Characterizing microservice dependency and performance: {Alibaba}
  trace analysis.
\newblock In \emph{Proceedings of ACM SoCC}, pages 412--426, 2021.

\bibitem{borgtraces2019}
M.~Tirmazi, A.~Barker, N.~Deng, M.~E. Haque, Z.~G. Qin, S.~Hand,
  M.~Harchol-Balter, and J.~Wilkes.
\newblock Borg: The next generation.
\newblock In \emph{Proceedings of EuroSys}, pages 30:1--30:14, 2020.

\bibitem{shahrad2020azure}
M.~Shahrad, R.~Fonseca, I.~Goiri, G.~Chaudhry, P.~Batum, J.~Cober,
  J.~Flores, and R.~Bianchini.
\newblock Serverless in the wild: Characterizing and optimizing the serverless
  workload at a large cloud provider.
\newblock In \emph{Proceedings of USENIX ATC}, pages 205--218, 2020.

\bibitem{semgrep2023}
Semgrep.
\newblock Semgrep: Lightweight static analysis for many languages.
\newblock \url{https://semgrep.dev/}, 2023.

\bibitem{cncfconformance}
Cloud Native Computing Foundation.
\newblock {CNCF} {Kubernetes} conformance.
\newblock \url{https://www.cncf.io/certification/software-conformance/}, 2023.

\bibitem{beyer2016sre}
B.~Beyer, C.~Jones, J.~Petoff, and N.~R. Murphy.
\newblock \emph{Site Reliability Engineering: How {Google} Runs Production
  Systems}.
\newblock O'Reilly Media, 2016.

\end{thebibliography}

\end{document}
